% concatenated from main.tex, chapters/1_introduction.tex, chapters/2_preliminaries.tex, chapters/3_problem_statement.tex, chapters/4_main_idea.tex, chapters/5_illustrative_example.tex, chapters/6_conclusions.tex
\documentclass[journal,twoside,web]{ieeecolor}
\usepackage{lcsys}
\usepackage{cite}
\usepackage{amsmath,amssymb,amsfonts}
\usepackage{algorithmic}
\usepackage{graphicx}
\usepackage{textcomp}
\def\BibTeX{{\rm B\kern-.05em{\sc i\kern-.025em b}\kern-.08em
    T\kern-.1667em\lower.7ex\hbox{E}\kern-.125emX}}
\markboth{\journalname, VOL. XX, NO. XX, XXXX 2017}
{Author \MakeLowercase{\textit{et al.}}: Preparation of Papers for textsc{IEEE Control Systems
Letters} (November 2021)}

\usepackage{array,multicol,multirow,booktabs,longtable,tabularx}
\usepackage{subfig}
\usepackage{amssymb,mathtools,dsfont,bbm,cuted,stmaryrd}
\usepackage{xspace,algorithmic,algorithm}


\usepackage{tikz}
\usetikzlibrary{math,arrows.meta,spy,automata,intersections,decorations,arrows,decorations.pathmorphing,decorations.markings,decorations.footprints,fadings,calc,trees,mindmap,shadows,decorations.text,patterns,positioning,shapes,matrix,fit,3d,plotmarks,svg.path}

\usepackage{pgfplots}


\usepgfplotslibrary{patchplots}

\usepackage{tikz-3dplot}
\usepackage{tkz-euclide}

\usepackage{placeins}
\usepackage{cuted}

% \usepackage{flushend}

\renewcommand{\algorithmicrequire}{\textbf{Input:}}
\renewcommand{\algorithmicensure}{\textbf{Output:}}

\let\labelindent\relax
\usepackage{enumitem}
\def \singlefigure {.7\textwidth}
\def \doublefigure {.45\textwidth}

\newcommand{\propref}[1]{{Proposition}~\ref{#1}\xspace}

\newcommand{\be}{\begin{equation}}
\newcommand{\ee}{\end{equation}}

\def\bal#1\eal{\begin{align}#1\end{align}}
\def\baln#1\ealn{\begin{align*}#1\end{align*}}

\newcommand{\ben}{\begin{equation*}}
\newcommand{\een}{\end{equation*}}

\newcommand{\re}[1]{\mathbb{R}^{#1}}%


\newcommand{\bbm}{\begin{bmatrix}}
\newcommand{\ebm}{\end{bmatrix}}

\newcommand{\bBm}{\begin{Bmatrix}}
\newcommand{\eBm}{\end{Bmatrix}}

\newcommand{\bvm}{\begin{vmatrix}}
\newcommand{\evm}{\end{vmatrix}}

\newcommand{\bVm}{\begin{Vmatrix}}
\newcommand{\eVm}{\end{Vmatrix}}

\newcommand{\bpm}{\begin{pmatrix}}
\newcommand{\epm}{\end{pmatrix}}

\newcommand{\bnm}{\begin{matrix}}
\newcommand{\enm}{\end{matrix}}

\newcommand{\bi}{\begin{itemize}}
\newcommand{\ei}{\end{itemize}}

\newcommand{\bse}{\begin{subequations}}
\newcommand{\ese}{\end{subequations}}

 %%%%%%%%%%%%%%%%%%%%%%%%%%%%%%%%%%%%%COMMON USED VARIABLES IN FTC%%%%%%%%%%%%%%%%%%%%%%%%%%%%%%%%%%%%%%%%%%%%%%%%%%%%
\newcommand{\xhat}[1]{\hat{x}_{#1}}
\newcommand{\xhatup}[1]{\hat{x_{#1}}^{UP}}
\newcommand{\xtilde}[1]{\tilde{x}_{#1}}

\newcommand{\zhat}[1]{\hat{z}_{#1}}
\newcommand{\zhatup}[1]{\hat{z}_{#1}^{UP}}


\newcommand{\I}{\mathcal{I}}

\newcommand{\sih}{\mathcal{I}_H}
\newcommand{\sif}{\mathcal{I}_F}
\newcommand{\sir}{\mathcal{I}_R}

\newcommand{\sihr}{\mathbf{I}_H}
\newcommand{\sifr}{\mathbf{I}_F}
\newcommand{\sirr}{\mathbf{I}_R}


\newcommand{\secref}[1]{Section~\ref{#1}\xspace}
\newcommand{\appref}[1]{Appendix~\ref{#1}\xspace}
\newcommand{\remref}[1]{Remark~\ref{#1}\xspace}
\newcommand{\lemref}[1]{Lemma~\ref{#1}\xspace}
\newcommand{\thmref}[1]{Theorem~\ref{#1}\xspace}
\newcommand{\cororef}[1]{Corollary~\ref{#1}\xspace}
\newcommand{\algref}[1]{Algorithm~\ref{#1}\xspace}
\newcommand{\assumref}[1]{Assumption~\ref{#1}\xspace}
\newcommand{\figref}[1]{Fig.~\ref{#1}\xspace}
\newcommand{\tabref}[1]{Table~\ref{#1}\xspace}
\newcommand{\dfref}[1]{Def.~\ref{#1}\xspace}

\newenvironment{proof-sketch}{\noindent{ \textit{Sketch of Proof}:}\hspace*{0.5em}}

\DeclareMathOperator{\diag}{diag}
%%%%%%%%%%%%%%%%%%%%%%%%%%%%%%%%%%%%%%%%%%%%%%%%%%%%%%%%%%%%%%%%%%%%%%%%%%%%%%%%%%%%%%%%%%%%%%%%%%%%%%%%

% \theoremstyle{plain}
\newtheorem{thm}{Theorem}
\newtheorem{prop}{Proposition}
\newtheorem{cor}{Corollary}
\newtheorem{lem}{Lemma}

%\theoremstyle{definition}
\newtheorem{df}{Definition}
%\theoremstyle{remark}
\newtheorem{rem}{Remark}
\newtheorem{assum}{Assumption}
\newtheorem{hyp}{Hypothesis}

\newcommand{\eot}{\ensuremath{\hfill\square}}
\newcommand{\eor}{\ensuremath{\hfill\blacklozenge}}
\def\qed{\hfill\vrule height 1.6ex width 1.5ex depth -.1ex}

\DeclareMathOperator{\vol}{Vol}
%\DeclareMathOperator{\det}{det}
\DeclareMathOperator{\dom}{dom}

% \title{Computing the Maximal Positive Invariant set\\ in the reduced space - zero pole\textcolor{red}{s} removal}
\title{On maximal positive invariant set computation for rank-deficient linear systems}
\author{Bogdan Gheorghe, Daniel Ioan, Cristian Flutur\\ Ionela Prodan, and Florin Stoican
\thanks{Bogdan Gheorghe, Daniel Ioan, Cristian Flutur, and Florin Stoican are with Faculty of Automation Control and Computer Science, National University of Science and Technology Politehnica Bucharest, Romania {\tt\small \{bogdan.gheorghe1807, daniel\char`_mihail.ioan, cristian.flutur, florin.stoican\}@upb.ro}}%
\thanks{Ionela Prodan is with the Univ. Grenoble Alpes, Grenoble INP$^\dagger$, LCIS, F-26000, Valence, France, $^\dagger$ Institute of Engineering and Management Univ. Grenoble Alpes.
        {\tt\small ionela.prodan@lcis.grenoble-inp.fr}}%
% %\thanks{Manuscript received March 8, 2024; revised May 10, 2024; accepted June 2, 2024. Date of publication June XX, 2024; date of current version June xx, 2019. The work of Bogdan Gheorghe was supported by a grant from the National Program for Research of the National Association of Technical Universities - GNAC ARUT 2023; Project ID: 207}
}

\usepackage{graphicx}      % include this line if your document contains figures
% \usepackage{natbib}        % required for bibliography
% \usepackage{enumitem}

% \usepackage{fontspec}

\usepackage{amsmath,amssymb,mathtools,stmaryrd}
% \usepackage{amsmath,amsfonts,amssymb}% amsthm
\usepackage{xcolor}


\usepackage[algo2e, linesnumbered, algoruled]{algorithm2e}

% Added by GB
\usepackage{hyperref}
\DeclareMathOperator*{\argmax}{arg\,max}
\DeclareMathOperator*{\argmin}{arg\,min}
\usepackage{svg}

% \usepackage{enumerate}

\iffalse
\newcommand{\bbm}{\begin{bmatrix}}
\newcommand{\ebm}{\end{bmatrix}}

\newtheorem{rem}{Remark}
\newtheorem{prop}{Proposition}
\newtheorem{coro}{Corollary}
\newtheorem{lem}{Lemma}

\newcommand{\eor}{\hfill\ensuremath{\blacklozenge}}
\fi

% \DeclareMathOperator{\det}{det}
% \DeclareMathOperator{\diag}{diag}
%===============================================================================
\pgfplotsset{compat=1.18}

\begin{document}


\pagestyle{empty} % do not insert page numbers
\maketitle
\thispagestyle{empty}


\begin{abstract}
The maximal positively invariant (MPI) set is obtained through a backward reachability procedure involving the iterative computation and intersection of predecessor sets under state and input constraints. 

However, standard static feedback synthesis may place some of the closed-loop eigenvalues at zero, leading to rank-deficient dynamics. This affects the MPI computation by inducing projections onto lower-dimensional subspaces during intermediate steps. By exploiting the Schur decomposition, we explicitly address this singular case and propose a robust algorithm that computes the MPI set in both polyhedral and constrained-zonotope representations.
\end{abstract}

\begin{IEEEkeywords}
maximal positive invariant set, constrained
zonotope, rank-deficient linear systems.
\end{IEEEkeywords}


%{\Large\textcolor{red}{Deadline CDC with LCSS: 17th of March}
%\textcolor{red}{Location CDC 2026: Hawaii, USA}}


\section{Introduction}
\label{sec:intro}
% \begin{itemize}
%     \item discussions MPI
%     \item the rest of the paper is organized as follow
%     \item tables notation
%     \item cite Kolmanovsky - LCSS
% \end{itemize}

The computation reachable sets is a fundamental problem at the intersection of control theory, algebraic geometry, and convex optimization, serving to define the exact theoretical bounds for safe system operation \cite{ossareh_complexity_2024,marpi2024lcss}. As a cornerstone of modern constrained control, reachability theory provides the mathematical guarantees necessary for recursive feasibility\cite{chen1998quasi}  and stability in Model Predictive Control (MPC)\cite{rakovic2018handbook} applications across the aerospace, chemical, and robotics sectors. %\cite{nguyen2020stability}.

A prime example of reachable set is the Maximal Positive Invariant (MPI) set which denotes the largest positive invariant set which also respects feasibility constraints. It is computed via a set recurrence relation involving the iterative calculation and intersection of predecessor sets subject to state and input constraints. Its primary use is a terminal set in Model Predictive Control\cite{rakovic2018handbook}. MPC operates by iteratively solving a finite-horizon optimization problem at each sampling interval, applying only the initial control action before receding the prediction horizon. A primary vulnerability is the potential loss of recursive feasibility: an optimal state sequence calculated at time $k$ may drive the plant to a state at $k+1$ from which future constraint violations become unavoidable. To prevent this and ensure closed-loop stability and constraint admissibility beyond the prediction horizon, the MPC formulation typically forces the terminal state into an MPI set where the control action switches to a fixed terminal feedback law.



% The necessity of the MPI set is particularly pronounced within MPC frameworks.
 

The feedback gain $K$ for this terminal law is usually computed via the resolution of a discrete-time algebraic Riccati equation. This unconstrained optimal design frequently results in placing one or more closed-loop eigenvalues at the origin ($0$). While mathematically optimal, this renders the resulting closed-loop state transition matrix rank-deficient, which presents significant geometric and computational challenges for calculating the required terminal MPI set. 
% Formally, the exact MPI set is computed via a set recurrence relation, involving the iterative calculation and intersection of predecessor sets subject to state and input constraints.
While the MPI remains full-dimensional it is the result of computations where the singular closed-loop matrix maps the system dynamics into a lower-dimensional subspace, fundamentally altering the structural properties of the intermediate sets generated during this recurrence sequence.

While established methodologies can compute these invariant sets for standard polyhedral constraints, current literature lacks efficient solutions for the severe computational burdens introduced by the set recurrence at high dimensions. To mitigate these, our recent work \cite{10552328} has leveraged constrained zonotopes. By adding algebraic equality constraints, these structures retain efficient linear operations while remaining closed under arbitrary intersections. This allows the exact MPI set recurrence to be executed entirely within the constrained zonotope framework. 
%Bypassing explicit half-space enumeration, this approach yields a succinct representation that drastically reduces computation time and memory.



%To address the  , this paper
Our idea is straightforward but, to the best of our knowledge, was never treated in the literature:
\begin{enumerate}[label=\roman*)]
    \item exploit the Schur factorization of the closed-loop matrix to identify the sub-space corresponding to non-zero eigenvalues and compute the reduced MPI set;
    \item lift the reduced MPI back into the original space and intersect it with ``initial constraints'' to output the MPI set w.r.t. the original dynamics and feasible set. 
\end{enumerate}
% We carry the computations for both polyhedral and constrained zonotopic representations. 

% \vspace{10em}


The remainder of this paper is organized as follows. \secref{sec:prel} provides the preliminary mathematical background on MPI sets and their geometrical representations. \secref{sec:justication} introduces the problem and provides theoretical justification for rank-deficient closed loop systems. \secref{sec:main_idea} details the MPI-computing procedure in both the polyhedral and constrained zonotopic representations. \secref{sec:ill_example} provides illustrative examples, while \secref{sec:conclusions} draws the conclusions.


\noindent \textit{Notations}: For a $x \in \mathbb R^{n}$, its infinity norm as $\|x\|_{\infty} := \max(|x_1|, |x_2|, \ldots, |x_n|)$. For $A \in \mathbb R^{n\times n}$, $\text{eig}(A)$ is  the set of all its eigenvalues. $I_n\in \mathbb R^{n\times n}$ is the identity matrix, and $\mathbf{1}_n$ is the column vector of $n$ values of one 

%$\mathbb N$ is the set of natural numbers. $O_{m \times n}\in \mathbb R^{m\times n}$ is the matrix with $m$ rows and $n$ columns whose entries are zero. Whenever $m=n$, we use the shorthand notation $O_{n}$. 

%The identity matrix is denoted by $I_n\in \mathbb R^{n\times n}$, and the symbol $\mathbf{1}_n$ represents the column vector of $n$ values of one. 

%For an arbitrary matrix $G\in \mathbb R^{m\times n}$, $G(i, \cdot) \in \mathbb R^{1\times n}$ its i-th row, and $G(\cdot, j)  \in \mathbb R^{m\times 1}$ its j-th column. $|A|$, with $A \in \mathbb R^{m\times n}$, represents the element-wise absolute value of $A$. For a vector $x \in \mathbb R^{n}$, its $1$ norm is given as $\|x\|_{1} := \sum_{i=1}^{i=n}|x_i|$ and its infinity norm as $\|x\|_{\infty} := \max(|x_1|, |x_2|, \ldots, |x_n|)$. 

%Any order operations (e.g., ``$\leq$'' or ``$\geq$'') or operator (e.g., ``$\log (\cdot)$'') are taken element-wise over the vector(s) considered. The function $\diag(v)$ constructs a square matrix, with the elements of the vector $v$ on its main diagonal and zeros elsewhere.

%The literature provides well-established methods for computing maximal (robust) positive invariant (M(R)PI) sets for linear systems with polyhedral constraints. However, a comprehensive understanding of the limitations and computational burdens associated with different methodologies is still lacking, most notably in the context of high-dimensional systems.

%the precise computation of the Maximal Positively Invariant (MPI) set represents a definitive, high-stakes intersection of control theory, advanced algebraic geometry, and convex optimization. By mathematically establishing the absolute exact bounds within which any dynamical system can safely operate indefinitely, the MPI set secures the fundamental architecture of modern constrained control. Specifically, it serves as the absolute linchpin for mathematically guaranteeing recursive feasibility and permanent stability in Model Predictive Control formulations across the aerospace, chemical, and robotics industries.

%The expansive mathematical theory of the Maximal Positively Invariant set achieves its greatest practical and commercial utility within the specific domain of Model Predictive Control (MPC). The MPC architecture functions by solving a massive finite-horizon optimal control optimization problem at every discrete sampling interval, applying only the very first calculated control action to the physical plant before immediately shifting the prediction horizon one step forward and recalculating. The overarching vulnerability of the MPC algorithm is the catastrophic loss of recursive feasibility; a mathematically optimal trajectory computed at time step $k$ may lead the physical plant to a state at time step $k+1$ from which absolutely no physically allowable control inputs can prevent a future, inevitable constraint violation.

%When dealing with discrete-time Linear Time-Invariant (LTI) systems, a "rank-deficient" or singular state transition matrix ($A$) introduces specific geometric and computational nuances to the calculation of the Maximal Positively Invariant (MPI) set. Rank deficiency means that the matrix maps the state space into a lower-dimensional subspace, effectively "compressing" the system's dynamics and creating a non-trivial null space.

%While convergence can be guaranteed, the rank deficiency alters the geometric properties of the intermediate sets during computation. For instance, when the matrix $A$ is singular, the successive sets in the sequence remain proper convex sets, but they may lose certain structural properties (such as being strictly semi-ellipsoidal) during the intermediate steps before finalizing into a proper convex polytope.


%A recent, highly significant mathematical breakthrough involves the deployment of constrained zonotopes. A constrained zonotope is a powerful extension that incorporates algebraic equality constraints on the generator coefficients. Constrained zonotopes brilliantly retain the highly efficient linear transformation and Minkowski sum properties of standard zonotopes, but crucially, because of the equality constraints, they are also perfectly closed under arbitrary intersections. This structural property means that the iterative intersection of input bounds and state bounds inherent in computing the exact MPI set can be performed strictly within the constrained zonotope mathematical framework. This entirely eliminates the need for the explicit enumeration of polyhedral half-spaces, providing a minimal, highly succinct, and exact representation of the MPI set that dramatically reduces both computational synthesis time and the necessary computer memory storage requirements.
\clearpage


\section{Preliminaries}
\label{sec:prel}

Consider the linear time-invariant system
\begin{equation}
\label{eq:dynamics_open_loop}
x_{k+1} = A x_k + B u_k,
\end{equation}
where $x_k, x_{k+1} \in \mathbb{R}^n$ denote the current and successor states; $u_k\in \mathbb{R}^m$ represents the input; $A \in \mathbb{R}^{n \times n}$, $B \in \mathbb{R}^{n \times m}$ are the state and input matrices. The sets $\mathcal{X} \subset \mathbb{R}^n$ and $\mathcal{U} \subset \mathbb{R}^m$ impose constraints on the system state and input, respectively.

% In order to ensure closed-loop stability and admissibility of the MPC algorithm, we need to gurantee that starting from $\bar x_{N}$ the state of the system will stay a MPI set $\Omega$ (i.e., $x_{N+\ell} \in \mathcal X, \bar u_{N+\ell}\in \mathcal U$ for any $\ell\geq 0$ if $x_N\in \Omega$).

Assuming the pair $(A,B)$ stabilizable, there exists a static gain $K \in \mathbb{R}^{m \times n}$ defining the feedback law $u_k = Kx_k$, which leads to the closed-loop matrix $A_\circ := A + BK$ that stabilizes the system. The resulting closed-loop dynamics of \eqref{eq:dynamics_open_loop} become
\begin{equation}
\label{eq:LTI_dynamics_closed}
x_{k+1} = A_\circ x_k .
\end{equation}

Furthermore, the open-loop constraints $x_k \in \mathcal{X}$ and $u_k \in \mathcal{U}$ imply the closed-loop state constraint $x_k \in \overline{\mathcal{X}}$ with 
\begin{equation}
\label{eq:barx}
\overline{\mathcal{X}} = \mathcal{X} \cap \{x:\, Kx \in \mathcal{U}\}.
\end{equation}

A popular tool in the arsenal of set-based methods is the so-called maximal positively invariant (MPI) set \cite{gilbert1991linear,kolmanovsky_theory_nodate}. It is often used to define the terminal set in model predictive control architectures \cite{rakovic2018handbook}, ensuring recursive feasibility since it is, by definition, safe (due to positive invariance) and non-conservative (the ``maximal'' property). For our purposes, it suffices to recall the standard set recurrence that computes it:
\begin{equation}
\label{eq:mpi_rec_standard}
\Omega_0=\overline{\mathcal X}, 
\qquad 
\Omega_{k+1}=A_\circ^{-1}\Omega_k \cap \overline{\mathcal X}.
\end{equation}
Under mild assumptions \cite[Thm.~3]{rakovic2022implicit}, the recurrence \eqref{eq:mpi_rec_standard} is guaranteed to terminate at a fixed point $\Phi_x := \Omega_{\bar k}=\Omega_{\bar k+1}$ for some finite index $\bar k$. Thus, an equivalent form to \eqref{eq:mpi_rec_standard} is
\begin{equation}
    \label{eq:mpi_rec_standard_2}
    \Phi_x=\bigcap\limits_{k=0}^{\bar k} \bigl\{x:\: A_\circ^kx \in \overline{\mathcal X}\bigr\}.
\end{equation}
Further insights and variations of the basic MPI construction can be found in \cite{blanchini2008set,houska_polyhedral_2025} and subsequent works.

\subsection*{Set descriptions}

While the recurrence \eqref{eq:mpi_rec_standard} is agnostic with respect to the nature of $\overline{\mathcal X}$, particular choices do influence the subsequent computations. A traditional option is the polyhedral representation \cite{ziegler2012lectures}. Writing $\mathcal X=\{F_Xx\leq \theta_X\}$, $ \mathcal U=\{F_Uu\leq \theta_U\}$ in their half-space form allows to express \eqref{eq:barx} as
\begin{equation}
\label{eq:hrep}
\overline{\mathcal X}:=\{F x \leq \theta\}=\biggl\{\bbm F_X\\ F_UK\ebm x\leq \bbm \theta_X\\\theta_U\ebm\biggr\},
\end{equation}
allows the recurrence \eqref{eq:mpi_rec_standard} to be written explicitly as
\begin{equation}
\label{eq:mpi-recurrence-polyhedral}
    \Phi_x=\bigcap\limits_{k=0}^{\bar k} \bigl\{x:\: FA_\circ^kx \leq \theta \bigr\}.
\end{equation}

An alternative representation is that of the constrained zonotope (CZ), first introduced in \cite{scott2016constrained} and increasingly used for set manipulations \cite{raghuraman2022set}. We call the compact set $\mathcal{CZ} \subset \re{n}$ a constrained zonotope if there exists $c \in \mathbb R^{n}, G \in \mathbb R^{n \times D}, F \in \mathbb \re^{p\times D}, \theta \in \mathbb R^{p}$ such that
\begin{equation}
\label{eq:con_poly_zono}
    \mkern-8mu\mathcal{CZ} = \bigl\langle c, G, F, \theta\bigr\rangle = \big \{x = c + G\lambda,\: F\lambda=\theta,\: \|\lambda\|_\infty \leq 1\bigr\}.
\end{equation}

% \begin{itemize}
%     \item polyhedral and constrained zonotopic descriptions \cite{houska_polyhedral_2025}
%     \item \textcolor{red}{see discussion about redundancy check} (sec 2.4, at the end): Polyhedral control design: Theory and methods
% \end{itemize}



% \noindent or using the convex hull of vertices (\textit{V-rep}):
% \be
% \label{eq:vrep} 
% P=\{x \in \re{d} : x=V\alpha, \alpha^\top \mathbf 1=1, \alpha \geq 0\}.
% \ee


% \begin{rem}
% A zonotope (Z) is a special case of a CZ with no equality constraints, i.e., $F$ and $\theta$ are empty. \eor
% \end{rem}

Noteworthy, CZs are closed under set intersection \cite{scott2016constrained}:
\begin{equation}
\label{eq:closedIntersection}
    \mkern-10mu \mathcal{CZ}_1 \cap \mathcal{CZ}_2 = \left\langle \mkern-4mu c_1, \bbm G_1 & 0 \ebm,\mkern-4mu  \bbm F_1 & \hphantom{-}0 \\ 0 & \hphantom{-}F_2 \\ G_1 & -G_2 \ebm, \bbm \theta_1 \\ \theta_2 \\ c_2 - c_1\ebm \right\rangle.
\end{equation}
Considering the state and input constraint sets as constrained zonotopes\footnote{In many practical cases the constraint sets are hyperrectangles, and therefore equality constraints are not required in their definition.}
$\mathcal X=\langle c_X, G_X, [\;], [\;]\rangle, \quad \mathcal U=\langle c_U, G_U, [\;], [\;]\rangle$, allows, via \eqref{eq:closedIntersection}, the set \eqref{eq:barx} to be expressed\footnote{We abuse the notation and re-use $F,\theta$ for the equality part of the CZ description.} as a CZ:
\begin{multline}
\label{eq:x_constr_zon}
\overline{\mathcal X} = \left\langle c, G, F, \theta\right\rangle =\bigl\langle c, \bbm G& 0\ebm,\\ \bbm KG& -G_U\ebm, c_U-Kc\bigr\rangle.
\end{multline}
Based on \cite{10552328}, all sets $\Omega_k = \left\langle c_k, G_k, F_k, \theta_k\right\rangle$ in the recurrence \eqref{eq:mpi_rec_standard} are also CZs and can be computed recursively as
\begin{multline}
\label{eq:mpi_constr_zon}
\Omega_{k+1} =\Bigl\langle A_\circ^{-1}c_k, \bbm A_\circ^{-1}G_k& 0\ebm,\\
    \bbm F_k&0\\ 0& F\\ A_\circ^{-1}G_k& -G\ebm,  \bbm\theta_k\\\theta\\c-A_\circ^{-1}c_k\ebm\Bigr\rangle.
\end{multline}

\vspace{-1.5em}
\subsection*{Illustrative example}
We take the discrete invariant linear dynamics defined by: \newline
$$A = \bbm 1.38 & 0.76 \\ 0.16 & 1.87 \ebm,\: B = \bbm 1 \\ 1\ebm$$
with input and state constraints given by $\mathcal X = \{\|x\|_{\infty} \leq 1\}$, $\mathcal U=\{\|u\|_{\infty} \leq 1\}$. Typically, closing the loop can be done either by solving the algebraic Riccati equation or by doing explicit pole placement. Each returns a stabilizing static gain $K$. Computing the MPI for dynamics \eqref{eq:LTI_dynamics_closed} and feasible set \eqref{eq:barx} for each gain $K$ results in the iterations illustrated in Figure~\ref{fig:example_LTI_MPI}. 
\begin{figure}[!ht]
\centering
\subfloat[MPI set for Ricatti]{\label{fig:example_LTI_idare}\includegraphics[width=0.5\columnwidth]{./pics/example_LTI_idare}}\hfill
\subfloat[MPI set for pole placement]{\label{fig:example_LTI_place}\includegraphics[width=0.5\columnwidth]{./pics/example_LTI_place}}
\caption{Example for computing the MPI set.}
\label{fig:example_LTI_MPI}
\end{figure}
Table~\ref{tab:MPI_method_comparison} shows that, even for this proof-of-concept example, significant differences appear in closed-loop in set shape and the value of $\bar k$, the index where the recurrence stops.%in the shape of the MPI and in the number of iterations needed to achieve it. 
\begin{table}[!ht]
\centering
\renewcommand{\arraystretch}{1.35}
\setlength\tabcolsep{3pt}

\begin{tabular}{c|c|c|c}
Method & Poles & Gain ($K$) & Nr. of iter. ($k$)\\\hline
Ricatti & $\bbm 0.48 & 0.83\ebm$ & $\bbm 2.73 & -0.80\ebm$ & 3\\\hline
Placement & $\bbm 0.95 & 0.7\ebm$ & $\bbm 1.43 & 0.16\ebm$ & 7
\end{tabular}

\caption{Comparison of Ricatti and pole placement.}
\label{tab:MPI_method_comparison}
\setlength\tabcolsep{6pt}
\end{table}


% We observed that, by changing the values of the poles for the pole placement method, the closer both poles are to one of the stability limits (i.e., $-1$ or $1$), the more iterations the MPI takes. This behavior is happening because \textcolor{red}{TBC by control guys}.

\clearpage


\section{Problem formulation and justification}
\label{sec:justication}
% \begin{itemize}
%     \item system theory: about poles
%     \item deflation
%     \item closed loop matrix, why we have zeros in closed loop
% \end{itemize}

% The closed-loop dynamics become
% \begin{equation}
%     \label{eq:cl-dynamics}
%     x^+ = A_\circ x
% \end{equation}

The state matrix of the closed-loop dynamics \eqref{eq:LTI_dynamics_closed} may possess multiple eigenvalues at $\lambda = 0$ (equivalently, $0 \in \mathrm{eig}(A_\circ)$), possibly with different partial multiplicities. Such situations are not exceptional; rather, they naturally arise in a variety of standard control synthesis procedures.
% 
% The state matrix of closed-loop dynamics \eqref{eq:LTI_dynamics_closed} may have more than one eigenvalue at $\lambda=0$ (equivalently put, $\{0\}\in \text{eig}(A_\circ)$, with different partial multiplicities. Such cases are not outliers but may actually arise from a variety of situations which are part of the normal control synthesis process. 

One such situation is optimal control, where the closed loop system may get poles placed anywhere in the unit disc, provided the stabilizing solution is computed \cite{ionescu1997general}. We recall that the stabilizing solution of the Riccati equation is computed as the disconjugated invariant subspace of the Hamiltonian matrix, therefore the closed loop gain matrix $K$ may place the eigenvalues of $A+BK$ anywhere in the stable region, so it may even place more than one eigenvalue at $\lambda=0$. Another situation of interest may result from a dead-beat control strategy \cite{van1984deadbeat}, where all closed loop eigenvalues are placed at $\lambda=0$.

\subsection*{Numerical example}

For illustration and later use in the paper, we consider the following numerical example, where the gain matrix $K$ assigns parts of the closed loop spectrum to $\lambda=0$. Consider the controllable pair $(A,B)$ and the feedback gain matrix $K$: 

\[
{\small
A = \begin{bmatrix}
-14.85 & -5.20 & -14.75 & -11.90 & -20.10 & -14.55 \\
-\hphantom{0}8.85 & \hphantom{-}0.10 & -12.95 & -\hphantom{0}9.20 & -10.20 & -13.15 \\
\hphantom{-0}9.90 & \hphantom{-}6.60 & \hphantom{-}10.30 & \hphantom{-0}6.80 & \hphantom{-}13.80 & \hphantom{-}10.10 \\
-14.95 & -7.50 & -13.85 & -10.20 & -21.00 & -13.65 \\
-18.40 & -5.70 & -26.40 & -17.70 & -23.10 & -26.40 \\
-12.35 & -3.80 & -21.85 & -13.30 & -14.90 & -21.85
\end{bmatrix}
}
\]
$$
B =
\begin{bmatrix}
1 & 4 \\
3 & 4 \\
0 & 0 \\
0 & 2 \\
4 & 4 \\
4 & 2
\end{bmatrix},
\quad K =
\begin{bmatrix}
1 & 0 & 4 & 2 & 1 & 4 \\
3 & 1 & 2 & 2 & 4 & 2
\end{bmatrix}.
$$
For the resulting closed loop matrix $A_\circ=A+BK$, we compute the Real Schur form matrix ($A_\circ = USU^\top$). Matrix
\[
S = \left[\begin{tabular}{ccc|cc|c}
0.70 & -0.87 & 0.49 & -0.31 & \hphantom{-}35.16 & -13.00 \\
\hphantom{0.0}0 & \hphantom{-}0.50 & 0.12 & -0.23 & \hphantom{-0}2.27 & \hphantom{0}-0.85 \\
\hphantom{0.0}0 & \hphantom{-0.0}0 & 0.20 & -0.40 & \hphantom{0}-3.46 & \hphantom{-0}2.33 \\
\hline
\hphantom{0.0}0 & \hphantom{-0.0}0 & \hphantom{0.0}0 & -0.00 & \hphantom{0}-2.14 & \hphantom{-0}1.25 \\
\hphantom{0.0}0 & \hphantom{-0.0}0 & \hphantom{0.0}0 & \hphantom{-}0.00 & \hphantom{0}-0.00 & \hphantom{0}-0.00 \\
\hline
\hphantom{0.0}0 & \hphantom{-0.0}0 & \hphantom{0.0}0 & \hphantom{-0.0}0 & \hphantom{-00.0}0 & \hphantom{0}-0.00
\end{tabular}\right]
\]
% \[
% U = \begin{bmatrix}
% \hphantom{-}0.07 & \hphantom{-}0.37 & -0.74 & -0.32 & 0.30 & -0.35 \\
% -0.34 & -0.52 & \hphantom{-}0.22 & -0.43 & 0.11 & -0.60 \\
% -0.67 & -0.16 & -0.20 & \hphantom{-}0.27 & 0.54 & \hphantom{-}0.35 \\
% \hphantom{-}0.61 & -0.65 & -0.24 & \hphantom{-}0.16 & 0.34 & \hphantom{-}0.09 \\
% \hphantom{-}0.13 & \hphantom{-}0.30 & \hphantom{-}0.35 & \hphantom{-}0.55 & 0.45 & -0.52 \\
% \hphantom{-}0.20 & \hphantom{-}0.23 & \hphantom{-}0.43 & -0.56 & 0.54 & \hphantom{-}0.35
% \end{bmatrix}
% \]
has its eigenvalues at $\{0,0,0,0.2,0.5,0.7\}$. Furthermore, the eigenvalues at $\lambda=0$ belong to two Jordan cells: one of dimension $1$ and one of dimension $2$, \cite[Chapter 3]{horn2012matrix}.

Such a simple example already reveals the issue: either explicitly, as in \eqref{eq:mpi-recurrence-polyhedral}, or implicitly, as in \eqref{eq:mpi_constr_zon}, we manipulate powers of the form
\[
A_\circ^k = (USU^\top)^k = U S^k U^\top.
\]
This becomes numerically unstable whenever $S$ is singular, as is the case here. Therefore, the remainder of the paper exploits the structure of $S$ to perform computations in a well-behaved subspace, before lifting the results back to the original dimension.






\newpage


\section{Main idea}
\label{sec:main_idea}
In the singular case, matrix $A_\circ$ is no longer full rank. Consequently, its image is not full dimensional. As a result, set recurrence \eqref{eq:mpi_rec_standard} may behave erratically or fail altogether, since multiplication by $A_\circ^k$ effectively acts as a projection onto the subspace associated with nonzero poles. In the remainder of this section, we propose constructions that explicitly handle poles at zero, regardless of their nature, for both polyhedral and constrained-zonotope set descriptions.

\subsection{The polyhedral case}

We handle first the case of polyhedral set descriptions \eqref{eq:hrep}--\eqref{eq:mpi-recurrence-polyhedral}, as shown in the following result.

\begin{prop}
\label{prop:p}
Consider the singular closed-loop matrix 
\begin{equation}
\label{eq:schur}
    A_\circ=A+BK = U\left[\begin{tabular}{c|c}
         $S_{11}$& $S_{12}$ \\\hline
         0& $S_{22}$
    \end{tabular}\right]U^\top, 
\end{equation}
whose Schur decomposition has the $(2,2)$ block $S_{22}\in \mathbb{R}^{d_2\times d_2}$ nilpotent of degree $p+1<n$, i.e.,
\begin{equation}
\label{eq:nilpotent}
    \bigl(S_{22}\bigr)^{p+1} = 0, \qquad \text{eig}(S_{22}) = 0.
\end{equation}
The scalars $d_1>0$, $d_2>0$ satisfy $d_1+d_2=n$.

Introducing the shorthand notation
\begin{equation}
    \label{eq:T}
    T=\sum_{i=0}^{p} \bigl(S_{11}\bigr)^{-i}S_{12}\bigl(S_{22}\bigr)^{i},
\end{equation}
the MPI set associated with the state matrix \eqref{eq:schur} and constraints \eqref{eq:hrep} is given by

\begin{equation}
    \label{eq:mpi-schur}
    \Phi_x = \Phi^{[1]}_x \cap \Phi^{[2]}_x, 
\end{equation}


\noindent with
% \begin{subequations}
\begin{align}
    \label{eq:mpi-schur_1}
    \Phi_x^{[1]} &= \bigcap\limits_{k=0}^{p-1} \overline{\mathcal X}=\bigcap\limits_{k=0}^{p-1}\biggl\{x:\:FA_\circ^k x\leq \theta\biggr\}, \\
    \label{eq:mpi-schur_2}
    \Phi_x^{[2]} &=\bigcap\limits_{k=p}^{\bar k +p}\overline{\mathcal X}= \biggl\{x:\: F_z \bbm \bigl(S_{11}\bigr)^p &\bigl(S_{11}\bigr)^{p-1}T\ebm U x\leq \theta_z\biggr\}.
    \end{align}
 % \end{subequations}
The pair $(F_z, \theta_z)$ defines 
\begin{equation}
        \label{eq:mpi_z}
    \Phi_z\hphantom{^{1}} =\bigcap\limits_{k=0}^{\bar k}\mathcal Z=\{z:\: F_zz\leq \theta_z\},
\end{equation}
the MPI set, computed as in \eqref{eq:mpi_rec_standard}, associated with the dynamics $z^+=S_{11}z$ and constraint set $\mathcal{Z}=\{z:\,F U_1 z\leq\theta\}$. $U_1$ gathers the first $d_1$ columns of $U$ and $U_2$ the remaining $d_2$ columns.
\end{prop}

% \textcolor{red}{U transpose?}

\begin{proof}
Under the Schur decomposition \eqref{eq:schur}, the dynamics \eqref{eq:LTI_dynamics_closed} and the associated set $\overline{\mathcal X}$ can be expressed in the transformed coordinates $y = U^\top x$, leading to
\begin{subequations}
\begin{align}
    \label{eq:y_dynamics}
    y_{k+1}&=\left[\begin{tabular}{c|c}
         $S_{11}$& $S_{12}$ \\\hline
         0& $S_{22}$
    \end{tabular}\right]y_k, \\
    \label{eq:y_set}\mathcal Y &= \{y:\: FUy\leq \theta\}.
\end{align}
\end{subequations}
Both the dynamics and the associated constraint set follow from the change of variables $x = Uy$ together with the orthogonality property $UU^\top = U^\top U = I$.

Decomposing $y=\bbm y_1^\top & y_2^\top \ebm^\top$ we use \eqref{eq:y_dynamics} to write the evolution of each component of the state
\begin{subequations}
\label{eq:dynamics}
    \begin{align}
        \nonumber y_{1, k+1} &= S_{11}\,y_{1, k} + S_{12}y_{2, k}\\
        &=\bigl(S_{11}\bigr)^{k+1}y_{1,0}+\sum\limits_{i=0}^k \bigl(S_{11}\bigr)^{k-i}S_{12}\bigl(S_{22}\bigr)^{i}y_{2,0},\\
        \nonumber y_{2, k+1} &= S_{22}y_{2, k}\\
        &=\bigl(S_{22}\bigr)^{k+1}y_{2,0}.
    \end{align}
\end{subequations}
% For further use we make the shorthand notation $T_{k,i}=\bigl(S_{11}\bigr)^{k-i}S_{12}\bigl(S_{22}\bigr)^{i}$. 
Exploiting the nilpotence condition \eqref{eq:nilpotent} we rewrite \eqref{eq:dynamics} into
\begin{subequations}
\label{eq:dynamics_nilpotent}
    \begin{align}
        \nonumber y_{1, k+1} =&  \bigl(S_{11}\bigr)^{k+1}y_{1,0}\\
        &+          \label{eq:dynamics_nilpotent_a}\bigl(S_{11}\bigr)^{k}\sum\limits_{i=0}^{\min(k,\, p)} \bigl(S_{11}\bigr)^{-i}S_{12}\bigl(S_{22}\bigr)^{i}y_{2,0},\\
        \label{eq:dynamics_nilpotent_b}y_{2, k+1} =& \begin{cases}\bigl(S_{22}\bigr)^{k+1}y_{2,0}, & k<p,\\
        0, & k\geq p.
        \end{cases}
    \end{align}
\end{subequations}
% Clearly, for any $k\geq p$, component $y_{2,k+1}=0$ as per \eqref{eq:dynamics_nilpotent_b}, and only $y_{1,k+1}$ updates as in \eqref{eq:dynamics_nilpotent_a}.
For convenience, let us denote $z_k=y_{1,k+p}$ and take $T$ as in \eqref{eq:T}. Consequently, for any $k\geq p$, \eqref{eq:dynamics_nilpotent_a} is equivalently written as (we shift from $k$ to $k-p$)
\begin{equation}
\label{eq:z}
    z_{k+1} = S_{11} z_k, \quad z_0=\bigl(S_{11}\bigr)^py_{1,0}+\bigl(S_{11}\bigr)^{p-1}Ty_{2,0}.
\end{equation}

Using \eqref{eq:dynamics_nilpotent} in \eqref{eq:y_set} we observe a similar behavior for the constraints: for $k\geq p$ only the first component ($y_{1, k}$) is affected: 
\begin{multline}
\label{eq:z_set}
    \mathcal Y=\{FUy_{k}\leq \theta\}=\bigl\{F\bbm U_1 & U_2\ebm \bbm y_{1, k}\\ y_{2, k}\ebm\leq \theta\bigr\}\\
    \overset{\eqref{eq:dynamics_nilpotent}}{=}\{FU_1 y_{1,k}\leq \theta\}\overset{\eqref{eq:z}}{=}\{FU_1 z_{k-p}\leq \theta\}. 
\end{multline}
At this point, we have both ingredients required for the set recurrence \eqref{eq:mpi_rec_standard}: the dynamics given by the first part of \eqref{eq:z} and the constraint set $\mathcal{Z}=\{z:\: FU_1 z\leq \theta\}$. Applying the recurrence yields \eqref{eq:mpi_z}. Next, rewriting $\Phi_z$ using the second part of \eqref{eq:z} allows us to constrain explicitly $y_{1,0}$ and $y_{2,0}$, obtaining \eqref{eq:mpi-schur_2} which enforces the constraints $y_k \in \mathcal Y$ for all $k \geq p$. By additionally imposing the constraints $y_k \in \mathcal Y$ for $k \leq p$, and reversing the change of coordinates $x = Uy$, we obtain \eqref{eq:mpi-schur_1}. Intersecting them gives \eqref{eq:mpi-schur}, thus concluding the proof.
\end{proof}


% \newpage 


% \newpage
\subsection{Constrained zonotopic case}

The same line of reasoning may be employed in the constrained zonotopic case \eqref{eq:mpi_constr_zon}--\eqref{eq:mpi_constr_zon}, as illustrated in the next result. 
\begin{prop}
\label{prop:cz}
    Consider the notation from \eqref{eq:schur}--\eqref{eq:T} and the constraint set given in its constrained zonotopic form $\overline{\mathcal X}=\langle c,G,F,\theta\rangle$, defined as in \eqref{eq:x_constr_zon}. Then the associated MPI set is given by 
    \begin{multline}
    \label{eq:mpi-schur-cz}
        \mkern-36mu\Phi_x = \mkern-2mu\biggl\langle \mkern-4mu c^{[1]},\mkern-2mu \bbm G^{[1]} & 0\ebm\mkern-2mu,\mkern-4mu \bbm F^{[1]} & 0\\ 0 & F^{[2]}\\ \bbm \bigl(S_{11}\bigr)^p &\bigl(S_{11}\bigr)^{p-1}T\ebm U^\top G^{[1]} & -G^{[2]}\mkern-4mu\ebm\mkern-4mu,\\
        \bbm \theta^{[1]}\\ \theta^{[2]}\\ c^{[2]}-\bbm \bigl(S_{11}\bigr)^p &\bigl(S_{11}\bigr)^{p-1}T\ebm c^{[1]}\ebm\biggr\rangle
    \end{multline}
    where the ``$[1]$'' parameters describe the constraints imposed until $k<p$:
    \begin{equation}
    \label{eq:omega_x_1_cz}
        \Phi_x^{[1]}=\bigcap\limits_{k=0}^{p-1} \overline{\mathcal X} = \bigl\langle c^{[1]}, G^{[1]}, F^{[1]}, \theta^{[1]}\bigr\rangle
    \end{equation}
and the ``$[2]$'' parameters describe the constraints applied to $k\geq p$ and, under dynamics \eqref{eq:z} and set recurrence \eqref{eq:mpi_constr_zon} which stops after $\bar k$ iterations \cite{10552328}:
    \begin{multline}
        \label{eq:omega_x_2_cz}
        \Phi_x^{[2]}=\bigcap\limits_{k=p}^{\bar k +p} {\mathcal X}\\
        =\{x:\: \bbm \bigl(S_{11}\bigr)^p &\bigl(S_{11}\bigr)^{p-1}T\ebm U^\top x\in \Phi_z\}, \theta^{[2]}\bigr\rangle.
    \end{multline}
    where
    \begin{equation}
        \label{eq:mpi_z_cz}
        \Phi_z =\bigcap\limits_{k=0}^{\bar k} {\mathcal Z}=\langle c^{[2]}, G^{[2]}, F^{[2]}, \theta^{[2]}\rangle,
    \end{equation}
    is the MPI set associated with the dynamics $z^+=S_{11}z$ and constraint set 
    \begin{equation}
    \label{eq:z-cz}
        \mathcal Z = \biggl\langle U_1^\top c, U_1^\top G, \bbm F\\ U_2^\top G\ebm, \bbm \theta\\ -U_2^\top c\ebm\biggr\rangle.
    \end{equation}
    % where, in \eqref{eq:schur}, $U= \bbm U_1 & U_2\ebm$, $U_1 \in \re{n\times d_1}$, $U_2 \in \re{n\times d_2}$. 
    All $(\cdot)^{[1]},(\cdot)^{[2]}$ parameters are obtained recursively as in \eqref{eq:mpi_constr_zon}.
\end{prop}
\begin{proof}
Using recurrence \eqref{eq:mpi_constr_zon} until step $k=p-1$ for closed-loop matrix \eqref{eq:schur} and constraint set \eqref{eq:x_constr_zon} gives $\Phi_x^{[1]}$ as in \eqref{eq:omega_x_1_cz}. Next, for $k\geq p$, we consider the same change of coordinates as before, $y=U^\top x$. Using definition \eqref{eq:con_poly_zono}, we have that $x_{k+1}\in \overline{\mathcal X}$ can be written as
\begin{subequations}
\label{eq:y_cz}
\begin{align}
\label{eq:y_cz_1}
    y_{1, k+1}&=U_1^\top c + U_1^\top G\lambda_{k+1},\quad F_1\lambda_{k+1} = \theta_1,\\
\label{eq:y_cz_2}
y_{2, k+1}&=U_2^\top c + U_2^\top G\lambda_{k+1},\quad F_2\lambda_{k+1} = \theta_2.
\end{align}    
\end{subequations}
Due to \eqref{eq:nilpotent} and as per \eqref{eq:dynamics_nilpotent_b} we have that $y_{2, k+1}=0$. This acts as an additional equality constraint on $\lambda_{k+1}$ in \eqref{eq:y_cz_2}. Introducing it into \eqref{eq:y_cz_1}, leads to inclusion $y_{1,k+1}\in \mathcal Z$, with $\mathcal Z$ defined as in \eqref{eq:z-cz}. With notation \eqref{eq:z} and the set \eqref{eq:z-cz}, set recurrence \eqref{eq:mpi_constr_zon} leads to \eqref{eq:mpi_z_cz}. Using the second part of \eqref{eq:z} and \eqref{eq:omega_x_1_cz} and \eqref{eq:omega_x_2_cz} leads to \eqref{eq:mpi-schur-cz}, thus concluding the proof.
\end{proof}

\subsection{Analysis}

Comparing Propositions~\ref{prop:p} and~\ref{prop:cz}, we observe that the same set names appear in both statements. This is not a mistake but a deliberate choice: the underlying geometric objects are identical. Indeed, constrained zonotopes and polyhedra have been shown to be equivalent set representations~\cite{raghuraman2022set}. The differences arise only in the specific algorithms used to manipulate them (in this case, the operations of matrix multiplication and set intersection). To emphasize this point, we present Algorithm~\eqref{alg:MPI_poles_in_zero}, which enumerates the computation steps for both polyhedral and constrained-zonotope representations.

\begin{algorithm}[!ht]
\SetKwComment{Comment}{}{}
\KwIn{$A_{\circ}$, $\overline{\mathcal{X}}$ as in \eqref{eq:hrep} or \eqref{eq:con_poly_zono}}
\KwOut{the MPI set $\Phi_x$}

\label{step:1}Compute the sorted real Schur form: $A_{\circ} = USU^T$ with $S$ and $p+1$ as in \eqref{eq:schur}
\BlankLine

% $\mathcal{Z}=\{z:\: FU_1 z\leq \theta\}$ computed as in \eqref{eq:z_set}
% \BlankLine

\textbf{Select procedure:\\ 
\hphantom{poly}
Polyhedral \hspace{2.56em} \vrule  \hspace{3em} Constrained Zonotope}
\BlankLine

$\mathcal{Z}$ from \eqref{eq:z_set} \; \hspace{4.25em} \vrule  \hspace{3.25em} $\,$ 
$\mathcal{Z}$ from \eqref{eq:z-cz}
\BlankLine

Take $A_\circ = S_{11}$, $\overline{\mathcal{X}}=\mathcal{Z}$ and compute reduced MPI set
\BlankLine

\hspace{3.77em}$\Phi_z$ as in \eqref{eq:mpi_z}\; \hspace{4.25em} \vrule   \hspace{3.5em} $\,$ $\Phi_z$ as in \eqref{eq:mpi_z_cz}
\BlankLine

Compute auxiliary sets:
\BlankLine

\hspace{3.5em} $\Phi_{x}^{[1]}$ from \eqref{eq:mpi-schur_1} \; \hspace{3.5em} \vrule   \hspace{3.5em} $\,$ 
$\Phi_{x}^{[1]}$ from \eqref{eq:omega_x_1_cz}
\BlankLine

\hspace{3.5em} Compute $T$ as in \eqref{eq:T}
\BlankLine

\hspace{3.5em} $\Phi_{x}^{[2]}$ from \eqref{eq:mpi-schur_2} \; \hspace{3.5em} \vrule   \hspace{3.5em} $\,$ 
$\Phi_{x}^{[2]}$ from \eqref{eq:omega_x_2_cz}
\BlankLine

Compute MPI set $\Phi_{x} = \Phi_{x}^{[1]} \bigcap \Phi_{x}^{[2]}$
\caption{MPI set for singular cases }
\label{alg:MPI_poles_in_zero}
\end{algorithm}
Noteworthy, step~1 of the algorithm implements the \emph{ordered Schur} decomposition from \cite{brandts2002matlab}, modified such that it arranges the Jordan and Schur blocks in descending order from largest to smallest, in magnitude. Thus, all the blocks with zero eigenvalue (if any) accumulate in the $(2,2)$ block, as shown in \eqref{eq:schur}.
% \newpage

\noindent Several remarks are in order. 
\begin{rem}
The set recurrence \eqref{eq:mpi_rec_standard} terminates at some finite index $\bar k$. Although unlikely, it is possible that $\bar k < p$. In this case, the construction in \eqref{eq:mpi-schur} or \eqref{eq:mpi-schur-cz} terminates without the need to compute \eqref{eq:mpi_z} or \eqref{eq:mpi_z_cz} or, indeed, the parts of \eqref{eq:mpi-schur_1}, \eqref{eq:omega_x_1_cz} which correspond to indices $\bar k+1, \ldots, p-1$. \eor
\end{rem}
\begin{rem}
Formulations \eqref{eq:mpi-schur} and \eqref{eq:mpi-schur-cz} also provide geometric insight into why a singular matrix disrupts the standard recurrence \eqref{eq:mpi_rec_standard}. The set $\Phi_z$ is computed in $\mathbb{R}^{d_1}$. When embedded into $\mathbb{R}^n$ as $\Phi_x^{[2]}$ after the change of variable $y=U^\top x$, it becomes unbounded along the subspace $\mathbb{R}^{d_2}$. Consequently, the bounds that ensure a well-behaved set must arise explicitly from the first $p$ iterations, as collected in \eqref{eq:mpi-schur_1} and \eqref{eq:omega_x_1_cz}, respectively.\eor
\end{rem}

% \begin{rem}
%     The maximal output admissible set (MOAS) is closely related \cite{gilbert1991linear} to the MPI set. Instead of constraining the initial state $x_0$ by the state's trajectory in the future ($x_k\in \overline{\mathcal X}, \forall k\geq 0)$ we constrain the initial state $x_0$ by the evolution of the output $\hat y_k = Cx_k$ in some output feasible domain ($\hat y_k\in \hat{\mathcal Y}$). Observing that the later inclusion is equivalent with checking $CA^k x_k \in \hat{\mathcal Y}, \forall k\geq 0$ it is straightforward to apply our setup by  $CA$  
% \end{rem}

\begin{rem}
Index $p+1$, as defined in \eqref{eq:nilpotent}, corresponds to the dimension of the largest Jordan cell of $S_{22}$. Inspecting the Jordan indices through the Schur form of $S_{22}$, \cite{kublanovskaya1982construction}, provides additional insight into the nature of the constraints that arise at a given time instant $k$. 

For illustration, suppose that the block dimensions belong to the (not necessarily unique) ordered set $\{p_1 \leq p_2 \leq \ldots \leq p_m\}$, with multiplicities $\{j_1, j_2, \ldots, j_m\}$. Then, for indices satisfying $p_i < k \leq p_{i+1}$, the constraints applied at step $k$ (i.e., $A^k x \in \overline{\mathcal X}$) describe a geometric object embedded in a subspace of dimension $n-\sum_{\ell=1}^{i} p_\ell j_\ell$. Once $k \geq p = \max_{\ell=1,\ldots,m} p_\ell$, the corresponding subspace dimension reaches $d_1 = n-d_2 = n-\sum_{\ell=1}^{m} p_\ell j_\ell$.
    % Index $p$, as defined in \eqref{eq:nilpotent} is only the largest Jordan block contained in $S_{22}$. Looking at all the Jordan and Schur blocks that describe $S_{22}$ offers additional insight in the nature of the constraints that apply at some time instant $k$. For the sake of argument, let us consider these dimensions to be in the possible not unique set $\{p_1\leq  p_2\leq \ldots, p_m\}$ having multiplicity $\{j_1, j_2, \ldots, j_m\}$. Thus, at step $p_i<k\leq p_{i+1}$, the constraints applied at this step ($A^kx\in \overline{\mathcal X}$) describe a geometric shape in the space of dimension $\sum\limits_{\ell=1}^i p_\ell j_\ell$. After $k\geq p = \max_{\ell=1,\ldots, m} = p_\ell$, the subspace dimension has become $d_1=\sum\limits_{\ell=1}^m p_\ell j_\ell$.
    \eor
\end{rem}

% Some remarks are in order. 

% % \begin{rem}
% % First step of the algorithm...\cite{}
% % \end{rem}




% % \newpage
% \begin{itemize}
% \item \textcolor{red}{Mention ordschur modification}
% \item algorithm for the conZ reduction
% % \item Minimal RPI, Maximal RPI, RPI
% \item illustrative example grafic: CSE the 2 methods of pole allocation
% \item \textcolor{red}{IMPLEMENT THE CONZ?}
% \item output constrained set? CA is not square but it can be completed with zeros and
% \item add here the example of 3 poles in zero @Bogdan
% \end{itemize}

% \newpage


\section{Illustrative example}
\label{sec:ill_example}
% \textcolor{red}{example from @Cristi, something more artificial - 1 column}

The following results were obtained using MPT3~\cite{MPT3} and CasADi~\cite{andersson2019casadi}. All computations were performed on a system running Ubuntu~22.04, equipped with an AMD Ryzen~9~3900X 3.8~GHz CPU (12 cores, 24 threads) and 48\,GB of RAM.

The MATLAB implementation corresponding to Prop.~\ref{prop:p} is available in the subdirectory {/replan-public/maximal-positive-invariant-set/LCSS-2026} of the GitLab repository:
\url{https://gitlab.com/replan/replan-public.git}.



\subsection{Numerical example in $\mathbb R^6$}

We recall the numerical example from \secref{sec:justication}. To this we attach state and input constraints $\mathcal X = \{x\in \mathbb R^{6}:\:\|x\|_{\infty} \leq 1\},\:  \mathcal U=\{u\in \mathbb R^{2}:\:\|u\|_{\infty} \leq 1\}$ which correspond to
\begin{equation*}
    \overline{\mathcal X}= \biggl\{\bbm I_6 \\ I_2K\ebm x\leq \bbm \mathbf{1}_6\\ \mathbf{1}_2\ebm\biggr\},
\end{equation*}
given as in \eqref{eq:barx}. Since in $S$, the upper-triangular matrix for the Schur factorization of $A_\circ$, we have two Jordan blocks (of dimensions $p_1=1$ and $p_2=2$), we have that the rank deficiency is $d_2 = p_1+p_2=3$ and $p=\max(p_1,p_2)=2$. Consequently, $d_1= n-d_2 = 3$ and $S_{11}$ is the $3\times 3$ matrix 
\begin{equation*}
    S_{11}=\begin{bmatrix}
0.70 & -0.87 & 0.49 \\
\hphantom{0.0}0 & \hphantom{-}0.50 & 0.12 \\
\hphantom{0.0}0 & \hphantom{-0.0}0 & 0.20
\end{bmatrix}.
\end{equation*}
Correspondingly, we compute the set $\mathcal Z\subset \mathbb R^3$ as shown below \eqref{eq:mpi_z}. Having both these elements we compute the reduced MPI set $\Phi_z$ from \eqref{eq:mpi_z} in $\overline{k}= 3$ iterations, as  depicted in \figref{fig:example_LTI_6dim}.


\begin{figure}[!ht]
\centering
\includegraphics[width=0.5\textwidth]{./pics/example_LTI_6dim}
% \includegraphics[width=\columnwidth]{./pics/example_LTI_6dim.svg}
\caption{MPI set $\Phi_z$ in $\mathbb R^3$.}
\label{fig:example_LTI_6dim}
\end{figure}

By selecting the blocks $S_{12}$ and $S_{22}$ from $S$ we compute, as in \eqref{eq:T}: 
\[
T = \begin{bmatrix}
-0.30 & 33.06 & -11.76 \\
-0.22 & \hphantom{0}2.23 & -\hphantom{0}0.82 \\
-0.39 & \hphantom{0}0.76 & -\hphantom{0}0.15
\end{bmatrix}
\]
which is used, as in \eqref{eq:mpi-schur_2} to obtain $\Phi_x^{[2]}\subset \mathbb R^6$. Set $\Phi_x^{[1]}\subset \mathbb R^6$ is obtained with $p=3$ iterations as in \eqref{eq:mpi-schur_1}. Intersecting the two, we obtain the desired result, the MPI set $\Phi_x$, characterizing the original dynamics and constraints. 
%The resulting MPI set $\Phi_{x} = \Phi_{x}^{[1]} \bigcap \Phi_{x}^{[2]}$ is depicted in \figref{fig:plot_MPI_numerical}.
Total computation time for all operations was $0.4$ seconds.

% \[
% S_{11} =
% \begin{bmatrix}
% 0.7000 & -0.8710 & 0.4931 \\
% 0      & 0.5000  & 0.1208 \\
% 0      & 0       & 0.2000
% \end{bmatrix}
% \]


\newpage

\subsection{The CSE example \cite{Leibfritz2006COMPleibCM}}
The Coupled Spring Experiment (CSE) considers a variable number of masses interconnected through springs and dampers. The system state consists of the positions and velocities of all masses, while the inputs are two external forces applied at the ends of the chain. The continuous-time model is given by
\begin{align}
\label{eq:cse1}
    \dot x &= 
    \underbrace{\begin{bmatrix}
    0 & I \\
    -M_{c}^{-1}K_{c} & -M_{c}^{-1}L_{c}
    \end{bmatrix}}_{\text{A}\in \mathbb R^{2\ell \times 2\ell}}x + 
    \underbrace{\begin{bmatrix}
    0 \\
    M_{c}^{-1}D_{c}
    \end{bmatrix}}_{\text{B}\in \mathbb R^{2\ell \times 2}} u,
\end{align}

where $M_{c} = \mu I$, $L_{c} = \delta I$, 
\begin{align}
    K_{c} = k\begin{bmatrix}
    1 & -1 & \cdots & 0 & 0 \\
    -1 & -2 & \ddots & 0 & 0 \\
    \vdots & \ddots & \ddots & \ddots & \vdots \\
    0 & 0 & \ddots & -2 & -1 \\
    0 & 0 & \cdots & -1 & 1
    \end{bmatrix} \nonumber, \; \text{and} \; D_{c} = \begin{bmatrix}
    1 & \hphantom{-}0 \\
    0 & \hphantom{-}0 \\
    \vdots & \hphantom{-}\vdots \\
    0 & \hphantom{-}0 \\
    0 & -1
    \end{bmatrix}, \nonumber
\end{align}
with $\mu = 4, \; \delta = 1, \; k = 1$. The system is discretized using forward Euler method with a sampling time of $1$ sec. The state and input constraints are: $\mathcal X = \{x\in \mathbb R^{2\ell}:\:\|x\|_{\infty} \leq 1\},\:  \mathcal U=\{u\in \mathbb R^{2}:\:\|u\|_{\infty} \leq 1\}$. By varying the number of masses through the parameter $\ell$, the state dimension becomes $2\ell$.
% Changing the number of masses in the system by changing the value of the parameter $\ell$, obtaining state dimension of $2\ell$.

The application of Algorithm~\ref{alg:MPI_poles_in_zero} in the polyhedral branch (Prop.~\ref{prop:p}) to the CSE dynamics, closed using both pole placement (all eigenvalues  in the unit disk) and a Riccati-based method (yielding a rank-deficient matrix), is illustrated in Fig.~\ref{fig:plot_data_simulations}. 

\begin{figure}[!ht]
\centering
\includegraphics[width=\columnwidth]{./pics/plot_data_simulations}
\caption{Plot number of iterations and commutation time.}
\label{fig:plot_data_simulations}
\end{figure}

We observe that:
\begin{itemize}
    \item Starting from higher dimensions ($\ell = 8$, corresponding to $x \in \mathbb{R}^{16}$), the number of iterations required by the standard algorithm ($\bar k$) increases significantly, whereas the increase is moderate for the rank-deficient method;
    \item The computation time further highlights the effectiveness of the proposed approach: the standard method becomes considerably more expensive for $\ell \geq 8$, while the rank-deficient method remains computationally efficient;
    \item The number of inequalities $\bar q$ required to represent the resulting polyhedral MPI sets is also larger for the standard approach (starting from $\ell = 5$). For instance, for $\ell = 10$, the standard method yields $\bar q = 1496$ inequalities, whereas the rank-deficient method reduces to $\bar q = 704$.
\end{itemize}
% \begin{itemize}
%     \item starting from higher dimensions ($\ell = 8$, corresponding to $x\in \mathbb R^{16}$), for the standard algorithm, the number of iterations needed ($\bar k$) increases considerable, while for the rank deficient method, the increase is modest;
%     \item the computation time emphasizes the effectiveness of the proposed method. The standard method takes much more time from $\ell \geq 8$, compared to the rank deficient method;
%     \item the number of inequalities $\bar q$, needed to represent the resulting polyhedral MPI sets, is also higher for the standard case (starting from $\ell = 5$). For example, computing the MPI for $\ell = 10$ with the standard method, results in a MPI set that has $\bar q = 1496$ inequalities, while for the rank deficient method the number drops to $\bar q = 704$.
% \end{itemize}

The resulting MPIs sets are not geometrically the same. This outcome is expected because of the pole placement, which changes the initial set $\overline{\mathcal{X}}$ (as shown in Fig.~\ref{fig:example_LTI_MPI}).

% \begin{figure}[!ht]
% \centering
% \includegraphics[width=\columnwidth]{./pics/plot_nr_ineq_simulation}
% \caption{Plot number of inequalities.}
% \label{fig:plot_nr_ineq_simulation}
% \end{figure}
% ~\newpage


\section{Conclusions}
\label{sec:conclusions}

% \vspace{12em}

We have presented a robust algorithm for computing the MPI set in the presence of rank-deficient dynamics, using the Schur factorization. The approach accommodates multiple zero eigenvalues with arbitrary multiplicities, and provides the resulting MPI set in both polyhedral and constrained-zonotope representations. Future work will extend this framework to the more general Maximal Output Admissible Set (MOAS), where we anticipate employing an SVD-based factorization.

% We have provided a robust algorithm to compute the MPI set in the case of rank-deficient matrices through the Schur factorization. We can handle multiple poles at zero with varying multiplicities (Jordan blocks of size greater than 1). The MPI is outputted in both polyhedral and constrained zonotope forms. 

% Future work will extended the work to the more general Maximal Output Admissible set (MOAS) where we expect to use the SVD factorization. 

\begin{thebibliography}{10}
\providecommand{\url}[1]{#1}
\csname url@samestyle\endcsname
\providecommand{\newblock}{\relax}
\providecommand{\bibinfo}[2]{#2}
\providecommand{\BIBentrySTDinterwordspacing}{\spaceskip=0pt\relax}
\providecommand{\BIBentryALTinterwordstretchfactor}{4}
\providecommand{\BIBentryALTinterwordspacing}{\spaceskip=\fontdimen2\font plus
\BIBentryALTinterwordstretchfactor\fontdimen3\font minus \fontdimen4\font\relax}
\providecommand{\BIBforeignlanguage}[2]{{%
\expandafter\ifx\csname l@#1\endcsname\relax
\typeout{** WARNING: IEEEtran.bst: No hyphenation pattern has been}%
\typeout{** loaded for the language `#1'. Using the pattern for}%
\typeout{** the default language instead.}%
\else
\language=\csname l@#1\endcsname
\fi
#2}}
\providecommand{\BIBdecl}{\relax}
\BIBdecl

\bibitem{ossareh_complexity_2024}
H.~R. Ossareh and I.~Kolmanovsky, ``\BIBforeignlanguage{en}{On {Complexity} {Bounds} for the {Maximal} {Admissible} {Set} of {Linear} {Time}-{Invariant} {Systems}},'' \emph{\BIBforeignlanguage{en}{IEEE Transactions on Automatic Control}}, pp. 1--8, 2024.

\bibitem{marpi2024lcss}
A.~Dey and S.~Bhasin, ``Computation of maximal admissible robust positive invariant sets for linear systems with parametric and additive uncertainties,'' \emph{IEEE Control Systems Letters}, vol.~8, pp. 1775--1780, 2024.

\bibitem{chen1998quasi}
H.~Chen and F.~Allgower, ``A quasi-infinite horizon nonlinear model predictive control scheme with guaranteed stability,'' \emph{Automatica}, vol.~34, no.~10, pp. 1205--1217, 1998.

\bibitem{rakovic2018handbook}
S.~V. Rakovi{\'c} and W.~S. Levine, \emph{Handbook of model predictive control}.\hskip 1em plus 0.5em minus 0.4em\relax Springer, 2018.

\bibitem{10552328}
B.~Gheorghe, D.-M. Ioan, F.~Stoican, and I.~Prodan, ``Computing the maximal positive invariant set for the constrained zonotopic case,'' \emph{IEEE Control Systems Letters}, vol.~8, 2024.

\bibitem{gilbert1991linear}
E.~Gilbert and K.~Tan, ``{Linear systems with state and control constraints: the theory and application of maximal output admissible sets},'' \emph{IEEE Tran. on Automatic Control}, vol.~36, no.~9, pp. 1008--1020, 1991.

\bibitem{kolmanovsky_theory_nodate}
I.~Kolmanovsky and E.~G. Gilbert, ``{Theory} and computation of disturbance invariant sets for discrete-time linear systems,'' \emph{Mathematical Problems in Engineering}, vol.~4, no.~4, pp. 317--363, 1998.

\bibitem{rakovic2022implicit}
S.~V. Rakovi{\'c} and S.~Zhang, ``The implicit maximal positively invariant set,'' \emph{IEEE Tran. on Automatic Control}, vol.~68, no.~8, pp. 4738--4753, 2023.

\bibitem{blanchini2008set}
F.~Blanchini and S.~Miani, \emph{Set-theoretic methods in control}.\hskip 1em plus 0.5em minus 0.4em\relax Springer, 2008, vol.~78.

\bibitem{houska_polyhedral_2025}
B.~Houska, M.~A. Müller, and M.~E. Villanueva, ``\BIBforeignlanguage{en}{Polyhedral control design: {Theory} and methods},'' \emph{\BIBforeignlanguage{en}{Annual Reviews in Control}}, vol.~60, p. 100992, 2025.

\bibitem{ziegler2012lectures}
G.~M. Ziegler, \emph{Lectures on polytopes}.\hskip 1em plus 0.5em minus 0.4em\relax Springer Science \& Business Media, 2012, vol. 152.

\bibitem{scott2016constrained}
J.~K. Scott, D.~M. Raimondo, G.~R. Marseglia, and R.~D. Braatz, ``Constrained zonotopes: A new tool for set-based estimation and fault detection,'' \emph{Automatica}, vol.~69, pp. 126--136, 2016.

\bibitem{raghuraman2022set}
V.~Raghuraman and J.~P. Koeln, ``Set operations and order reductions for constrained zonotopes,'' \emph{Automatica}, vol. 139, p. 110204, 2022.

\bibitem{ionescu1997general}
V.~Ionescu, C.~Oara, and M.~Weiss, ``General matrix pencil techniques for the solution of algebraic riccati equations: a unified approach,'' \emph{IEEE Transactions on Automatic Control}, vol.~42, no.~8, pp. 1085--1097, 1997.

\bibitem{van1984deadbeat}
P.~Van~Dooren, ``Deadbeat control: A special inverse eigenvalue problem,'' \emph{BIT Numerical Mathematics}, vol.~24, no.~4, pp. 681--699, 1984.

\bibitem{horn2012matrix}
R.~A. Horn and C.~R. Johnson, \emph{Matrix analysis}.\hskip 1em plus 0.5em minus 0.4em\relax Cambridge university press, 2012.

\bibitem{brandts2002matlab}
J.~H. Brandts, ``Matlab code for sorting real {Schur} forms,'' \emph{Numerical linear algebra with applications}, vol.~9, no.~3, pp. 249--261, 2002.

\bibitem{kublanovskaya1982construction}
V.~Kublanovskaya, ``Construction of a canonic basis for matrices and pencils of matrices,'' \emph{Journal of Soviet Mathematics}, vol.~20, no.~2, pp. 1929--1942, 1982.

\bibitem{MPT3}
M.~Herceg, M.~Kvasnica, C.~Jones, and M.~Morari, ``{Multi-Parametric Toolbox 3.0},'' in \emph{Proc.~of the European Control Conference}, Z\"urich, Switzerland, July 17--19 2013, pp. 502--510.

\bibitem{andersson2019casadi}
J.~A. Andersson, J.~Gillis, G.~Horn, J.~B. Rawlings, and M.~Diehl, ``{CasADi}: a software framework for nonlinear optimization and optimal control,'' \emph{Mathematical Programming Computation}, vol.~11, no.~1, pp. 1--36, 2019.

\bibitem{Leibfritz2006COMPleibCM}
F.~Leibfritz, ``{COMPleib: COnstrained Matrix–optimization Problem library} – a collection of test examples for nonlinear semidefinite programs, control system design and related problems.''\hskip 1em plus 0.5em minus 0.4em\relax Department of Mathematics, Univ. Trier, Germany, Tech. Rep., 2006.

\end{thebibliography}



\end{document}
